\section{Introduction, motivation, etc.}

\lecture{1} For most of this lecture, we will be motivating why we care about
Lie algebras and why they are defined the way they are. If there are sections
of it that don't make sense, they can be skipped (although if none of it makes
sense, you should be worried).

Connected Lie groups are groups which are also differentiable manifolds. The
simple algebraic groups over algebraically closed fields (which form a nice
class of Lie groups) are:
\begin{align*}
  SL(n) &= \setc{A \in \Mat_n}{\det A = 1},\\
  SO(n) &= \setc{A \in SL_n}{AA^T = I},\\
  Sp(2n) &= \setc{A \in SL_{2n}}{A \text{ preserves a symplectic form}},
\end{align*}
and exactly $5$ others (as we will actually see!). (We are ignoring finite
centres, so $\Spin(n)$, the double cover of $SO(n)$, is missing.)

\begin{remark}
Each of these groups has a maximal compact subgroup. For example, $SU(n)$ is
the maximal compact subgroup of $SL_n(\C)$, while $SO_n(\R)$ is the maximal
compact subgroup of $SO_n(\C)$. These are themselves Lie groups, of course.
The representations of the maximal compact subgroup are the same as the
algebraic representations of the simple algebraic group, which in turn are the
same as the finite-dimensional representations of its (the simple algebraic
group's) Lie algebra. These last are what we study.
\end{remark}

Lie groups are complicated objects: for example, $SU_2(\C)$ is homeomorphic to
a $3$-sphere, and $SU_n$ is a twisted product of a $3$-sphere, a $5$-sphere,
\dots, a $(2n-1)$-sphere. Thus, studying Lie groups requires some amount of
algebraic topology and a lot of algebraic geometry. We want to replace these
``complicated'' disciplines by the ``easy'' discipline of linear algebra.

Therefore, instead of a Lie group $G$ we will be considering $\g = T_IG$, the
tangent space to $G$ at the identity.

\begin{definition}
A \emph{linear algebraic group} is a subgroup of some $GL_n$ defined by
polynomial equations in the matrix coefficients.
\end{definition}

We see that the examples above are linear algebraic groups.

\begin{remark}
It is not altogether obvious that this is an intrinsic characterisation (it
would be nice not to depend on the embedding into $GL_n$). The intrinsic
characterisation is as ``affine algebraic groups'', that is, groups which are
affine algebraic varieties and for which multiplication and inverse are
morphisms of affine algebraic varieties. One direction of this identification
is relatively easy (we just need to check that multiplication and inverse
really are morphisms); in the other direction, we need to show that any affine
algebraic group has a faithful finite-dimensional representation, i.e.\ embeds
in $GL(V)$. This involves looking at the ring of functions and doing something
with it.
\end{remark}

We will now explain a tangent space if you haven't met with it before.

\begin{example}
Let's pretend we're physicists, since we're in their building anyway. Let
$G = SL_2$, and let $|\eps| \ll 1$. Then for
\[
  g = \begin{pmatrix} 1 & \\ & 1 \end{pmatrix}
      + \eps \begin{pmatrix} a & b \\ c & d \end{pmatrix}
      + \text{higher-order terms}
\]
to lie in $SL_2$, we must have $\det g = 1$, or
\[
  1 = \det \begin{pmatrix} 1 + \eps a & \eps b \\ \eps c & 1 + \eps d \end{pmatrix}
      + \text{h.o.t.}
    = 1 + \eps(a + d) + \eps^2(\text{junk}) + \text{h.o.t.}
\]
Therefore, to have $g \in G$ we need to have $a + d = 0$.
\end{example}

To do this formally, we define the \emph{dual numbers}
\[
  E = \C[\eps]/\eps^2 = \setc{a + b\eps}{a, b \in \C}.
\]
For a linear algebraic group $G$, we consider
\[
  G(E) = \setc{A \in \Mat_n(E)}{A \text{ satisfies the polynomial equations defining } G \subset GL_n}.
\]
For example,
\[
  SL_2(E) = \Bigl\{ \begin{pmatrix} \alpha & \beta \\ \gamma & \delta \end{pmatrix}
     \Bigm| \alpha, \beta, \gamma, \delta \in E,\ \alpha\delta - \beta\gamma = 1 \Bigr\}.
\]
The natural map $E \to \C$, $\eps \mapsto 0$ gives a projection
$\pi\colon G(E) \to G$. We define the Lie algebra of $G$ as
\[
  \g \cong \pi^{-1}(I) \cong \setc{X \in \Mat_n(\C)}{I + \eps X \in G(E)}.
\]
In particular,
\[
  \sl_2 = \Bigl\{ \begin{pmatrix} a & b \\ c & d \end{pmatrix} \in \Mat_2(\C)
     \Bigm| a + d = 0 \Bigr\}.
\]

\begin{remark}
$I + \eps X$ represents an ``infinitesimal change at $I$ in the direction
$X$''. Equivalently, it is a map $\Spec \C[\eps]/\eps^2 \to G$, the
first-order germ of a curve through $I$.
\end{remark}

\begin{exercise}
(Do this only if you already know what a tangent space and tangent bundle
are.) Show that $G(E) = TG$ is the tangent bundle to $G$, and $\g = T_IG$ is
the tangent space to $G$ at $I$.
\end{exercise}

\begin{example}
Let $G = GL_n = \setc{A \in \Mat_n}{A^{-1} \text{ exists}}$.
Claim:
\begin{align*}
  G(E) &:= \setc{\tilde A \in \Mat_n(E)}{\tilde A^{-1} \text{ exists}}\\
       &= \setc{A + B\eps}{A, B \in \Mat_n(\C),\ A^{-1} \text{ exists}}.
\end{align*}
Indeed, $(A + B\eps)(A^{-1} - A^{-1}BA^{-1}\eps) = I$.
\end{example}

\begin{remark}
It is not obvious that $GL_n$ is itself a linear algebraic group (what are the
polynomial equations for ``determinant is non-zero''?). However, we can think
of $GL_n \subset \Mat_{n+1}$ as
\[
  \Bigl\{ \begin{pmatrix} A & \\ & \lambda \end{pmatrix}
     \Bigm| \det(A) \cdot \lambda = 1 \Bigr\}.
\]
\end{remark}

\begin{example}
$G = SL_n(\C)$.
\end{example}

\begin{exercise}
$\det(I + \eps X) = 1 + \eps \tr(X)$.
\end{exercise}

As a corollary to the exercise, $\sl_n = \setc{X \in \Mat_n}{\tr(X) = 0}$.

\begin{example}
$G = O_n(\C) = \setc{A}{AA^T = I}$. Then
\begin{align*}
  \g &:= \setc{X \in \Mat_n(\C)}{(I + \eps X)(I + \eps X)^T = I}\\
     &= \setc{X \in \Mat_n(\C)}{I + \eps(X + X^T) = I}\\
     &= \setc{X \in \Mat_n(\C)}{X + X^T = 0},
\end{align*}
the antisymmetric matrices. Now, if $X + X^T = 0$ then $2\tr(X) = 0$, and
since we're working over $\C$ and not in characteristic $2$, we conclude
$\tr(X) = 0$. Therefore, this is also the Lie algebra of $SO_n$.

This is not terribly surprising, because topologically $O_n$ has two connected
components corresponding to determinant $+1$ and $-1$ (they are images of each
other via a reflection). Since $\g = T_IG$, it cannot see the $\det = -1$
component, so this is expected.
\end{example}

The above example prompts the question: what exactly is it in the structure of
$\g$ that we get from $G$ being a group?

The first thing to note is that $\g$ does \emph{not} get a multiplication.
Indeed, $(I + A\eps)(I + B\eps) = I + (A + B)\eps$, which has nothing to do
with multiplication.

The operation that turns out to generalise nicely is the commutator,
$(P, Q) \mapsto PQP^{-1}Q^{-1}$. Taken as a map $G \times G \to G$ that sends
$(I, I) \mapsto I$, this should give a map $T_IG \times T_IG \to T_IG$ by
differentiation.

\begin{remark}
Generally, differentiation gives a \emph{linear} map, whereas what we will get
is a \emph{bilinear} map. This is because we will in fact differentiate each
coordinate: i.e.\ first differentiate the map $f_P\colon G \to G$,
$Q \mapsto PQP^{-1}Q^{-1}$ with respect to $Q$ to get a map
$\tilde f_P\colon \g \to \g$ (with $P \in G$ still) and then differentiate
again to get a map $\g \times \g \to \g$.
\end{remark}

What is this map? Let $P = I + A\eps$, $Q = I + B\delta$ where
$\eps^2 = \delta^2 = 0$ but $\eps\delta \neq 0$. Then
\begin{align*}
  PQP^{-1}Q^{-1} &= (I + A\eps)(I + B\delta)(I - A\eps)(I - B\delta)\\
                 &= I + (AB - BA)\eps\delta.
\end{align*}
Thus, the binary operation on $\g$ should be $[A, B] = AB - BA$.

\begin{definition}
The \emph{bracket} of $A$ and $B$ is $[A, B] = AB - BA$.
\end{definition}

\begin{exercise}
Show that $(PQP^{-1}Q^{-1})^{-1} = QPQ^{-1}P^{-1}$ implies $[A, B] = -[B, A]$,
i.e.\ the bracket is skew-symmetric.
\end{exercise}

\begin{exercise}
Show that the associativity of multiplication in $G$ implies the \emph{Jacobi
identity}
\[
  0 = [[X, Y], Z] + [[Y, Z], X] + [[Z, X], Y].
\]
\end{exercise}

\begin{remark}
The meaning of ``implies'' in the above exercises is as follows: we want to
think of the bracket as the derivative of the commutator, not as the explicit
formula $[A, B] = AB - BA$ (which makes the skew-symmetry obvious, and the
Jacobi identity only slightly less so). For example, we could have started
working in a different category.
\end{remark}

We will now define our object of study.

\begin{definition}
Let $K$ be a field, $\charac K \neq 2, 3$. A \emph{Lie algebra} $\g$ is a
vector space over $K$ equipped with a bilinear map (\emph{Lie bracket})
$[\,,\,]\colon \g \times \g \to \g$ with the following properties:
\begin{enumerate}
  \item $[X, Y] = -[Y, X]$ (skew-symmetry);
  \item $[[X, Y], Z] + [[Y, Z], X] + [[Z, X], Y] = 0$ (Jacobi identity).
\end{enumerate}
\end{definition}

\lecture{2} Examples of the Lie algebra definition from last time:

\begin{example}\leavevmode
\begin{enumerate}
  \item $\gl_n = \Mat_n$ with $[A, B] = AB - BA$.
  \item $\so_n = \setc{A}{A + A^T = 0}$.
  \item $\sl_n = \setc{A}{\tr(A) = 0}$ (note that while $AB$ need not have
    trace $0$ for $A, B \in \sl_n$, we do have $\tr(AB) = \tr(BA)$, so
    $[A, B] \in \sl_n$ even though $AB \notin \sl_n$ in general).
  \item $\spl_{2n} = \setc{A \in \gl_{2n}}{JA^TJ^{-1} + A = 0}$, where $J$ is
    the symplectic form
    \[
      J = \begin{pmatrix}
            &&&&& 1\\
            &&&& \iddots &\\
            &&& 1 &&\\
            && -1 &&&\\
            & \iddots &&&&\\
            -1 &&&&&
          \end{pmatrix}.
    \]
  \item $\fb = \Biggl\{ \begin{pmatrix}
            * & * & \cdots & *\\
              & * & \cdots & *\\
              &   & \ddots & \vdots\\
              &   &        & *
          \end{pmatrix} \in \gl_n \Biggr\}$,
    the upper triangular matrices (the name is $\fb$ for Borel, but we don't
    need to worry about that yet).
  \item $\fn = \Biggl\{ \begin{pmatrix}
            0 & * & \cdots & *\\
              & 0 & \cdots & *\\
              &   & \ddots & \vdots\\
              &   &        & 0
          \end{pmatrix} \in \gl_n \Biggr\}$,
    the strictly upper triangular matrices (the name is $\fn$ for nilpotent,
    but we don't need to worry about that yet).
  \item For any vector space $V$, let $[\,,\,]\colon V \times V \to V$ be the
    zero map. This is the ``abelian Lie algebra''.
\end{enumerate}
\end{example}

\begin{exercise}\leavevmode
\begin{enumerate}
  \item Check directly that $\gl_n$ is a Lie algebra.
  \item Check that the other examples above are Lie subalgebras of $\gl_n$,
    that is, vector subspaces closed under $[\,,\,]$.
\end{enumerate}
\end{exercise}

\begin{example}
$\begin{pmatrix} * & * \\ * & 0 \end{pmatrix}$ is \emph{not} a Lie subalgebra
of $\gl_2$.
\end{example}

\begin{exercise}
Find the algebraic groups whose Lie algebras are given above.
\end{exercise}

\begin{exercise}
Classify all Lie algebras of dimension $3$. (You might want to start with
dimension $2$. I'll do dimension $1$ for you on the board: skew-symmetry of
the bracket means that the bracket is zero, so there is only the abelian one.)
\end{exercise}

\begin{definition}
A \emph{representation} of a Lie algebra $\g$ on a vector space $V$ is a
homomorphism of Lie algebras $\phi\colon \g \to \gl_V$. We say that $\g$
\emph{acts on} $V$.
\end{definition}

Observe that, above, our Lie algebras were defined with a (faithful)
representation in tow.

There is always a representation coming from the Lie algebra acting on itself
(it is the derivative of the group acting on itself by conjugation):

\begin{definition}
For $x \in \g$ define $\ad(x)\colon \g \to \g$ as $y \mapsto [x, y]$.
\end{definition}

\begin{lemma}
$\ad\colon \g \to \End \g$ is a representation (called the \emph{adjoint
representation}).
\end{lemma}

\begin{proof}
We need to check whether $\ad([x, y]) = \ad(x)\ad(y) - \ad(y)\ad(x)$, i.e.\
whether the equality holds when we apply it to a third element $z$. But
\[
  \ad([x, y])(z) = [[x, y], z]
\]
and
\[
  \ad(x)\ad(y)(z) - \ad(y)\ad(x)(z) = [x, [y, z]] - [y, [x, z]]
    = -[[y, z], x] - [[z, x], y].
\]
These are equal by the Jacobi identity.
\end{proof}

\begin{definition}
The \emph{centre} of $\g$ is
$\setc{x \in \g}{[x, y] = 0 \ \forall y \in \g} = \ker(\ad\colon \g \to \End \g)$.
\end{definition}

Thus, if $\g$ has no centre, then $\ad$ is an embedding of $\g$ into
$\End \g$ (and conversely, of course).

Is it true that every finite-dimensional Lie algebra embeds into some $\gl_V$,
i.e.\ is linear? (Equivalently, is it true that every finite-dimensional Lie
algebra has a faithful finite-dimensional representation?) We see that if $\g$
has no centre, then the adjoint representation makes this true. On the other
hand, if we look inside $\gl_2$, then
$\fn = \bigl(\begin{smallmatrix} 0 & * \\ 0 & 0 \end{smallmatrix}\bigr)$ is
abelian, so $\ad$ maps it to zero in $\gl(\fn)$, despite the fact that we
started with $\fn \subset \gl_2$. That is, we can't always just take $\ad$.

\begin{theorem}[Ado's theorem; we will not prove it]
Any finite-dimensional Lie algebra over $K$ is a subalgebra of some $\gl_n$,
i.e.\ admits a faithful finite-dimensional representation.
\end{theorem}

\begin{remark}
This is the Lie algebra equivalent of the statement we made last time about
algebraic groups embedding into $GL(V)$. That theorem was ``easy'' because
there was a natural representation to look at (functions on the algebraic
group). There isn't such a natural object for Lie algebras, so Ado's theorem
is actually hard.
\end{remark}

\begin{example}
$\g = \sl_2$ with basis
\[
  e = \begin{pmatrix} 0 & 1 \\ 0 & 0 \end{pmatrix}, \quad
  h = \begin{pmatrix} 1 & 0 \\ 0 & -1 \end{pmatrix}, \quad
  f = \begin{pmatrix} 0 & 0 \\ 1 & 0 \end{pmatrix}
\]
(these are the standard names). Then $[e, f] = h$, $[h, e] = 2e$, and
$[h, f] = -2f$.
\end{example}

\begin{exercise}
Check this!
\end{exercise}

A representation of $\sl_2$ is a triple $E, F, H \in \Mat_n$ satisfying the
same bracket relations. Where might we find such a thing?

In this lecture and the next, we will get them as derivatives of
representations of $SL_2$. We will then rederive them from just the linear
algebra.

\begin{definition}
If $G$ is an algebraic group, an \emph{algebraic representation} of $G$ on $V$
is a homomorphism of groups $\rho\colon G \to GL(V)$ defined by polynomial
equations in the matrix coefficients. (This can be done so that it is
independent of the embedding into $GL_n$. Think about it.)
\end{definition}

To get a representation of the Lie algebra out of this, we again use the dual
numbers $E$. If we use $E = K[\eps]/\eps^2$ instead of $K$, we will get a
homomorphism of groups $G(E) \to GL_V(E)$. Moreover, since $\rho(I) = I$ and
it commutes with the projection map, we get
\[
  \rho(I + A\eps) = I + \eps \cdot (\text{some function of } A) = I + \eps\, d\rho(A)
\]
(this is to be taken as the definition of $d\rho(A)$).

\begin{exercise}
$d\rho$ is the derivative of $\rho$ evaluated at $I$ (i.e.\
$d\rho\colon T_IG \to T_IGL_V$).
\end{exercise}

\begin{exercise}
The fact that $\rho\colon G \to GL_V$ is a group homomorphism means that
$d\rho\colon \g \to \gl_V$ is a Lie algebra homomorphism, i.e.\ $V$ is a
representation of $\g$.
\end{exercise}

\begin{example}
$G = SL_2$. Let $L(n)$ be the space of homogeneous polynomials of degree $n$
in two variables $x, y$, with basis $x^n, x^{n-1}y, \dots, xy^{n-1}, y^n$ (so
$\dim L(n) = n + 1$). Then $GL_2$ acts on $L(n)$ by change of coordinates: for
$g = \bigl(\begin{smallmatrix} a & b \\ c & d \end{smallmatrix}\bigr)$ and
$f \in L(n)$ we have
\[
  \bigl((\rho_n g) f\bigr)(x, y) = f(ax + cy, bx + dy).
\]
In particular, $\rho_0$ is the trivial representation of $GL_2$, $\rho_1$ is
the usual $2$-dimensional representation $K^2$, and
\[
  \rho_2 \begin{pmatrix} a & b \\ c & d \end{pmatrix}
  = \begin{pmatrix} a^2 & ab & b^2 \\ 2ac & ad + bc & 2bd \\ c^2 & cd & d^2 \end{pmatrix}.
\]
Since $GL_2$ acts on $L(n)$, we see that $SL_2$ acts on $L(n)$.
\end{example}

\begin{remark}
The proper way to think of this is as follows. $GL_2$ acts on $\PP^1$ and on
$\mathcal{O}(n)$ on $\PP^1$, hence on the global sections
$\Gamma(\PP^1, \mathcal{O}(n)) = S^nK^2$. That's the source of these
representations. We can do this sort of thing for all the other algebraic
groups we listed previously, using flag varieties instead of $\PP^1$ and
possibly higher cohomology instead of the global sections (this is a theorem
of Borel, Weil, and Bott). That, however, requires algebraic geometry, and gets
harder to do in infinitely many dimensions.
\end{remark}

Differentiating the above representations, e.g.\ for
$e = \bigl(\begin{smallmatrix} 0 & 1 \\ 0 & 0 \end{smallmatrix}\bigr)$, we see
\[
  \rho(I + \eps e)\, x^iy^j = x^i(\eps x + y)^j = x^iy^j + \eps j x^{i+1}y^{j-1}.
\]
Therefore, $d\rho(e)\, x^iy^j = j x^{i+1}y^{j-1}$.

\begin{exercise}\leavevmode
\begin{enumerate}
  \item In the action of the Lie algebra,
    \begin{align*}
      e(x^iy^j) &= j x^{i+1}y^{j-1},\\
      f(x^iy^j) &= i x^{i-1}y^{j+1},\\
      h(x^iy^j) &= (i - j) x^iy^j.
    \end{align*}
  \item Check directly that these formulae give a representation of $\sl_2$.
  \item Check that $L(2)$ is the adjoint representation.
  \item Show that the formulae $e = x\frac{\partial}{\partial y}$,
    $f = y\frac{\partial}{\partial x}$,
    $h = x\frac{\partial}{\partial x} - y\frac{\partial}{\partial y}$ give an
    (infinite-dimensional!) representation of $\sl_2$ on $k[x, y]$.
  \item Let $\charac(k) = 0$. Show that $L(n)$ is an irreducible representation
    of $\sl_2$, hence of $SL_2$.
\end{enumerate}
\end{exercise}

\lecture{3} Last time we defined a functor
\begin{multline*}
  \bigl(\text{algebraic representations of a linear algebraic group } G\bigr)\\
  \longrightarrow
  \bigl(\text{Lie algebra representations of } \g = \Lie(G)\bigr)
\end{multline*}
via $\rho \mapsto d\rho$.

This functor is not as nice as you would like to believe, except for the Lie
algebras we care about.

\begin{example}
$G = \C^\times \implies \g = \Lie(G) = \C$ with $[x, y] = 0$.

A representation of $\g = \C$ on $V$ is a matrix $A \in \End(V)$ (to
$\rho\colon \C \to \End(V)$ corresponds $A = \rho(1)$).

$W \subseteq V$ is a submodule iff $AW \subseteq W$, and $\rho$ is isomorphic
to $\rho'\colon \g \to \End(V')$ iff $A$ and $A'$ are conjugate as matrices.

Therefore, representations of $\g$ correspond to the Jordan normal forms of
matrices.

As any linear transformation over $\C$ has an eigenvector, there's always a
one-dimensional subrepresentation of $V$. Therefore, $V$ is irreducible iff
$\dim V = 1$. Also, $V$ is \emph{completely decomposable} (a direct sum of
irreducible representations) iff $A$ is diagonalisable.

For example, if
\[
  A = \begin{pmatrix}
        0 & 1 &        &        &   \\
          & 0 & 1      &        &   \\
          &   & \ddots & \ddots &   \\
          &   &        & \ddots & 1 \\
          &   &        &        & 0
      \end{pmatrix}
\]
then the associated representation is indecomposable, but not irreducible.
The invariant subspaces are $\langle e_1\rangle, \langle e_1, e_2\rangle,
\dots, \langle e_1, e_2, \dots, e_n\rangle$, and none of the proper nonzero
ones has an invariant complement.
\end{example}

Now let's look at the representations of $G = \C^\times$.

The irreducible representations are $\rho_n\colon G \to GL_1 = \Aut(\C)$,
for $n \in \Z$, where $\rho_n(z)$ is multiplication by $z^n$. Every
finite-dimensional representation is a direct sum of these.

\begin{remark}
Proving these statements about representations of $G$ takes some theory of
algebraic groups. However, you probably know this for $S^1$, which is homotopy
equivalent to $G$. Recall that to get complete reducibility in the finite group
setting you took an average over the group; you can do the same thing here,
because $S^1$ is compact and so has finite volume. In fact, for the theory
that we study, our linear algebraic groups have maximal compact subgroups,
which have the same representations, which for this reason are completely
reducible. We will not prove that representations of the linear algebraic
groups are the same as the representations of their maximal compact
subgroups.
\end{remark}

Observe that the representations of $G$ and of $\g$ are not the same. Indeed,
$\rho \mapsto d\rho$ will send $\rho_n \mapsto n \in \C$ (prove it!).
Therefore, irreducible representations of $G$ embed into irreducible
representations of $\g$, but the map is not remotely surjective (not to
mention the decomposability issue).

The above example isn't very surprising, since $\g$ is also the Lie algebra
$\Lie(\C, +)$, and we should not expect representations of $\g$ to coincide
with those of $\C^\times$. The surprising result is

\begin{theorem}[Lie]
$\rho \mapsto d\rho$ is an equivalence of categories $\Rep G \to \Rep \g$ if
$G$ is a simply connected simple algebraic group.
\end{theorem}

\begin{remark}
A ``simple algebraic group'' is not simple in the usual sense: e.g.\ $SL_n$
has a centre, which is obviously a normal subgroup. However, if $G$ is simply
connected and simple in the above sense, then the centre of $G$ is finite
(e.g.\ $Z_{SL_n} = \mu_n$, the $n$th roots of unity), and the proper normal
subgroups are subgroups of the centre.
\end{remark}

\begin{exercise}
If $G$ is an algebraic group, and $Z$ is a finite central subgroup of $G$,
then $\Lie(G/Z) = \Lie(G)$. (Morally, $Z$ identifies points ``far away'', and
therefore does not affect the tangent space at $I$.)
\end{exercise}

We have now seen that the map (algebraic groups) $\to$ (Lie algebras) is not
injective (see the above exercise, or, more shockingly, $G = \C$ and
$G = \C^\times$). In fact,

\begin{exercise}
Let $G_n = \C^\times \ltimes \C$ where $\C^\times$ acts on $\C$ via
$t \cdot \lambda = t^n\lambda$ (so
$(t, \lambda)(t', \lambda') = (tt', \lambda + t^n\lambda')$). Show that
$G_n \cong G_m$ iff $n = \pm m$. Show that $\Lie(G_n) \cong \C x + \C y$ with
$[x, y] = y$ for every $n \neq 0$.
\end{exercise}

The map also isn't surjective (its image consists of the ``algebraic Lie
algebras''). This is easily seen in characteristic $p$; for example,
$\sl_p/\text{centre}$ cannot be the image of an algebraic group. In general,
algebraic groups have a Jordan decomposition (every element can be written as
$(\text{semisimple}) \times (\text{unipotent})$), and therefore the algebraic
Lie algebras should have a Jordan decomposition as well. In Bourbaki, you can
find an example of a $5$-dimensional Lie algebra for which the semisimple and
the nilpotent parts of its elements lie only in the ambient $\gl_V$ and not in
the algebra itself.
